\documentclass[12pt]{article}
\usepackage{amsmath,amssymb,amsfonts}
\usepackage[margin=1in]{geometry}
\begin{document}
\begin{center}
{\large\bfseries 13th Irish Mathematical Olympiad\\6 May 2000}
\end{center}
\bigskip\noindent\textbf{Paper 1}\bigskip
\begin{enumerate}
\item Let $S$ be the set of all numbers of the form $a(n) = n^2 + n + 1$, where $n$ is a natural number. Prove that the product $a(n)a(n+1)$ is in $S$ for all natural numbers $n$. Give, with proof, an example of a pair of elements $s, t \in S$ such that $st \notin S$.

\item Let $ABCDE$ be a regular pentagon with its sides of length one. Let $F$ be the midpoint of $AB$ and let $G$, $H$ be points on the sides $CD$ and $DE$, respectively, such that $\angle GFD = \angle HFD = 30^\circ$. Prove that the triangle $GFH$ is equilateral. A square is inscribed in the triangle $GFH$ with one side of the square along $GH$. Prove that $FG$ has length
\[
t = \frac{2\cos 18^\circ (\cos 36^\circ)^2}{\cos 6^\circ},
\]
and that the square has sides of length
\[
\frac{t\sqrt{3}}{2 + \sqrt{3}}.
\]

\item Let $f(x) = 5x^{13} + 13x^5 + 9ax$. Find the least positive integer $a$ such that 65 divides $f(x)$ for every integer $x$.

\item Let
\[
a_1 < a_2 < a_3 < \cdots < a_M
\]
be real numbers. $\{a_1, a_2, \ldots, a_M\}$ is called a \textbf{weak arithmetic progression} of length $M$ if there exist real numbers $x_0, x_1, x_2, \ldots, x_M$ and $d$ such that
\[
x_0 \leq a_1 < x_1 \leq a_2 < x_2 \leq a_3 < x_3 \leq \cdots \leq a_M < x_M
\]
and for $i = 0, 1, 2, \ldots, M-1$, $x_{i+1} - x_i = d$ i.e.\ $\{x_0, x_1, x_2, \ldots, x_M\}$ is an arithmetic progression.
\begin{enumerate}
\item[(a)] Prove that if $a_1 < a_2 < a_3$, then $\{a_1, a_2, a_3\}$ is a weak arithmetic progression of length 3.
\item[(b)] Let $A$ be a subset of $\{0, 1, 2, \ldots, 999\}$ with at least 730 members. Prove that $A$ contains a weak arithmetic progression of length 10.
\end{enumerate}

\item Consider all parabolas of the form $y = x^2 + 2px + q$ ($p, q$ real) which intersect the $x$- and $y$-axes in three distinct points. For such a pair $p$, $q$ let $C_{p,q}$ be the circle through the points of intersection of the parabola $y = x^2 + 2px + q$ with the axes. Prove that all the circles $C_{p,q}$ have a point in common.

\end{enumerate}
\newpage
\begin{center}
{\large\bfseries 13th Irish Mathematical Olympiad\\6 May 2000}
\end{center}
\bigskip\noindent\textbf{Paper 2}\bigskip
\begin{enumerate}
\setcounter{enumi}{5}
\item Let $x \geq 0$, $y \geq 0$ be real numbers with $x + y = 2$. Prove that
\[
x^2 y^2 (x^2 + y^2) \leq 2.
\]

\item Let $ABCD$ be a cyclic quadrilateral and $R$ the radius of the circumcircle. Let $a, b, c, d$ be the lengths of the sides of $ABCD$ and $Q$ its area. Prove that
\[
R^2 = \frac{(ab+cd)(ac+bd)(ad+bc)}{16Q^2}.
\]
Deduce that
\[
R \geq \frac{(abcd)^{3/4}}{Q\sqrt{2}},
\]
with equality \textbf{if and only if} $ABCD$ is a square.

\item For each positive integer $n$ determine with proof all positive integers $m$ such that there exist positive integers $x_1 < x_2 < \cdots < x_n$ with
\[
\frac{1}{x_1} + \frac{2}{x_2} + \frac{3}{x_3} + \cdots + \frac{n}{x_n} = m.
\]

\item Prove that in each set of ten consecutive integers there is one which is coprime with each of the other integers.
For example, taking 114, 115, 116, 117, 118, 119, 120, 121, 122, 123 the numbers 119 and 121 are each coprime with all the others. [Two integers $a$, $b$ are coprime if their greatest common divisor is one.]

\item Let $p(x) = a_0 + a_1 x + \cdots + a_n x^n$ be a polynomial with non-negative real coefficients. Suppose that $p(4) = 2$ and that $p(16) = 8$. Prove that $p(8) \leq 4$ and find, with proof, all such polynomials with $p(8) = 4$.

\end{enumerate}
\end{document}
